% ---------------------------------------------------------------
% seccoes/caracterizacao.tex — Aspetos caracterizadores
% ---------------------------------------------------------------

\chapter{Aspetos caracterizadores}\label{sec:2}

Identifica-se neste capítulo as características do sistema pericial implementado, assim como se apresenta as principais limitações tidas em conta no decorrer do respetivo desenvolvimento.

O que distingue o presente projeto de um sistema pericial de consulta é o domínio dinâmico, ou seja, o sistema desenvolvido não faz perguntas ao utilizador, uma vez que recebe observações dez vezes por segundo e tem de decidir antes que a situação mude. Daqui resultam os aspetos que caracterizam a solução apresentada:

\begin{itemize}
  \item Raciocínio temporal: as regras analisam janelas de eventos recentes para distinguir uma leitura atípica de um perigo persistente;
  \item Risco \emph{vs} prioridade: o nível de risco descreve a gravidade da situação, enquanto a prioridade, definida por categorias, decide entre propostas. Desta forma, a decisão final não depende da ordem pela qual as regras disparam;
  \item Dois canais de decisão conjunta: a velocidade e a via decidem em separado. Um veículo pode, no mesmo ciclo, abrandar e iniciar uma ultrapassagem;
  \item Mundo, perceção e conhecimento separados: o mundo simulado contém tudo o que existe, mas o motor de inferência só conhece o que os sensores detetaram. A título de exemplo, um peão na estrada não influencia a decisão enquanto não for visto;
  \item Utilização de estratégias como a lógica difusa para adaptar a velocidade à meteorologia;
  \item Explicabilidade completa: cada decisão regista a regra, a evidência e as propostas concorrentes;
  \item Tempo real observável: a interface recebe o estado completo do sistema a cada ciclo, por \emph{WebSocket}
\end{itemize}
